\uuid{JEUG}
\exo7id{819}
\titre{exo7 819}
\auteur{ridde}
\organisation{exo7}
\datecreate{1999-11-01}
\isIndication{false}
\isCorrection{false}
\chapitre{Calcul d'intégrales}
\sousChapitre{Intégration par parties}
\cpge{mpsi-mp2i.integration.c, pcsi.integration.b, ptsi.integration.b, mpsi-mp2i.nombres-reels-suites.d, pcsi.nombres-reels-suites.c, ptsi.nombres-reels-suites.c, mpsi-mp2i.procedes-sommatoires-discrets.d, bcpst1.derivees-primitives-integrales.d, bcpst1.suites-reelles, bcpst1.integration-segment.d, mp-mpi.series-numeriques-vectorielles.a, pt.series-numeriques.b, bcpst2.series-reelles}
\module{Analyse}
\niveau{L1}
\difficulte{}

\contenu{
\texte{
Soit $I_{n} = \int_0^1 \frac{x^n}{1 + x}dx$.
}
\begin{enumerate}
    \item \question{En majorant la fonction int\'egr\'ee, montrer que $ (I_{n})_{n \in \Nn} \rightarrow 0$.}
    \item \question{Calculer $I_n + I_{n + 1}$.}
    \item \question{D\'eterminer $\lim \limits_{n \rightarrow  + \infty} (\sum\limits_{k = 1}^n \frac{
 (-1)^{k + 1}}k)$.}
\end{enumerate}
}
